\begin{proof}

\begin{enumerate}[label={$\langle1\rangle$\arabic*.},leftmargin=*]
\item It suffices to check a critical value $c$. Represent a stalk class
by $a\in H^1(V;\Z)$ for $V=f^{-1}(\Delta_c)$ and a sufficiently
small disc $\Delta_c$. Put $V^*=f^{-1}(\Delta_c\setminus\{c\})$.
Assume that the restriction of $a$ to a smooth fibre is zero.
The homotopy sequence for $V^*\to\Delta_c^*$ shows that
$a|_{V^*}$ is the pullback of an integer multiple $k$ of the
positive meridian class of the punctured disc.

\item Choose a smooth point of one irreducible component of the central
divisor, away from all its other components. There the map has
the form $t=z^m u$ with $m>0$ and $u$ a holomorphic unit.
A small positive normal circle has base winding $m$ and bounds
its normal disc in $V$. Evaluating $a$ on that circle gives
$km=0$, whence $k=0$.
The inclusion $V^*\hookrightarrow V$ induces a surjection of
fundamental groups: a loop can be perturbed off a complex divisor,
including its higher-codimension singular strata. Therefore
$a|_{V^*}=0$ implies $a=0$.
Taking the direct limit over the discs proves the asserted stalk
injection; proper topological base change identifies its source
with $H^1(f^{-1}(c);\Z)$.

\item If a global section vanished at one point, it would vanish on a
neighbourhood of that point. Its restriction to the connected
smooth locus is a section of a local system, so it would vanish
everywhere on that locus. The just-proved injection then makes
it zero everywhere, as required.

\end{enumerate}
\end{proof}
